% Cómo funciona HTTPS: de la URL al candado — guía de estudio (skill /dossier). Qué pasa entre escribir una dirección https y ver el candado, y qué protege ese candado.
% Perfil: breve — texto normal — imágenes normal — lector estudio — formato digital — temas=4
% Pedido: «/dossier breve lector=estudio: un resumen para estudiar cómo funciona HTTPS: de la URL al candado (DNS, TCP, el handshake de TLS 1.3, certificados y cadena de confianza, qué protege y qué no)»
% Medir:  python3 ~/.claude/skills/dossier/scripts/medir.py breve-https.tex --paginas 5 --paginas-min 2 --paginas-texto 2.0 --parte-texto 0.4 --densa 550 --visuales-min 4 --lector estudio
% Regenerar:  python3 graficos.py && latexmk -lualatex breve-https.tex
\documentclass[10pt,a4paper]{article}
\usepackage[spanish,es-noshorthands]{babel}
\usepackage{dossier}
\documento{HTTPS · de la URL al candado}{Cómo funciona HTTPS: de la URL al candado}{Ejemplo de /dossier}

\begin{document}

% ============================ PÁGINA 1: todo el documento en una página ============================
\portada{Redes · guía de estudio}
  {Cómo funciona HTTPS: de la URL al candado}
  {Qué pasa entre escribir una dirección y ver el candado, y qué garantiza ese candado}
  {Resumen para estudiar. Cada dato lleva en gris su fuente; las siglas remiten a las referencias del final.}

\begin{cifras}[4]
  \cifra{Puerto por defecto}{443}{si la URL https no indica otro \fuente{RFC 9110, §4.2.2}}
  \cifra{Handshake de TLS 1.3}{1 RTT}{TLS 1.2 necesitaba 2 \fuente{RFC 8446 · RFC 5246}}
  \cifra{Hasta la respuesta}{3 RTT}{TCP, TLS 1.3 y HTTP, sin el DNS (cuenta propia)}
  \cifra{Controles por certificado}{4}{firma, vigencia, revocación y emisor \fuente{RFC 5280, §6.1.3}}
\end{cifras}

\resumen
\textbf{Del nombre a la conexión.} El navegador traduce el nombre a una dirección IP con
DNS, abre una conexión TCP al puerto 443 y sobre ella negocia TLS \fuente{RFC 9110, §4.3.3}.
TCP gasta un viaje de ida y vuelta (RTT) antes de que el cliente pueda enviar datos
\fuente{RFC 9293, §3.5}.

\textbf{El handshake.} En TLS 1.3, cliente y servidor acuerdan las claves en un solo RTT, y
todo lo que sigue al ServerHello viaja cifrado \fuente{RFC 8446, §1.2 y §2}. En una conexión
nueva, la primera respuesta HTTP llega a los 3 RTT.

\textbf{La identidad.} El certificado X.509 ata el nombre del servidor a una clave pública, y
el servidor firma el handshake con la clave privada. El navegador sigue la cadena de firmas
hasta una raíz que ya tiene \fuente{RFC 5280, §6.1}.

\textbf{El alcance.} HTTPS cifra la ruta, los encabezados y el cuerpo, y detecta a quien se
haga pasar por el servidor. No oculta la IP, el nombre del sitio ni el tamaño y el ritmo del
tráfico \fuente{RFC 8446, §1.2 y §E.3 · RFC 9110, §4.3.3}.

\ojo[Lo que el candado no dice]{Que el sitio sea honesto. Para emitir un certificado, Let's
Encrypt solo le pide a quien lo solicita que pruebe que controla el dominio
\fuente{LE-validación}.}

\begin{bloque}[El camino prepara el handshake; handshake y certificado determinan lo que HTTPS protege. Cada caja lleva a su sección.]
\begin{tikzpicture}
  \foreach \x/\capa in {0/EL CAMINO, 6.52/EL MECANISMO, 13.04/EL ALCANCE}
    \node[etiqueta,font=\scriptsize\ttfamily,text=rotulo,anchor=south west] at (\x,1.4) {\capa};
  \node[cajag,anchor=west] (m1) at (0,0)
    {\arriba{0.65cm}{3.9cm}{\hyperref[sec:camino]{\textbf{01 De la URL a TCP}}\\DNS y un RTT de TCP}};
  \node[cajag,anchor=west] (m2) at (6.52,0.65)
    {\arriba{0.65cm}{3.9cm}{\hyperref[sec:handshake]{\textbf{02 El handshake}}\\claves en un solo RTT}};
  \node[cajag,anchor=west] (m3) at (6.52,-0.65)
    {\arriba{0.65cm}{3.9cm}{\hyperref[sec:certificados]{\textbf{03 Certificado y cadena}}\\quién es el servidor}};
  \node[cajag,anchor=west] (m4) at (13.04,0)
    {\arriba{0.65cm}{3.9cm}{\hyperref[sec:alcance]{\textbf{04 Qué protege y qué no}}\\contenido sí, metadatos no}};
  \draw[bus] (m2.west) -- ++(-0.25,0) |- (m3.west);
  \draw[bus] (m2.east) -- ++(0.25,0) |- (m3.east);
  \draw[rel] (m1.east) -- node[relacion,above] {prepara} ($(m2.west)+(-0.25,-0.65)$);
  \draw[rel] ($(m2.east)+(0.25,-0.65)$) -- node[relacion,above] {determinan} (m4.west);
\end{tikzpicture}
\end{bloque}

% ============================ 01 ============================
\section{La URL fija el destino; DNS y TCP abren el camino}\label{sec:camino}

El esquema https obliga a asegurar la conexión antes de enviar la petición: el servidor
tiene que estar autenticado y la comunicación, cifrada y protegida contra cambios
\fuente{RFC 9110, §4.2.2}.

\begin{bloque}[Cada parte de la URL la usa una capa distinta. En azul, lo único que viaja cifrado, dentro de la petición HTTP \fuente{MDN-TLS}; el host viaja en claro ({\ver[sección 04]{sec:alcance}}).]
\begin{tikzpicture}
  \path (0,0.4) (17.4,0);
  \begin{scope}[every node/.style={anchor=base west,inner sep=1.5pt,font=\ttfamily\large,text=tinta}]
    \node (esq) at (5.0,0) {https};
    \node[text=apagado] (sep) at (esq.base east) {://};
    \node (host) at (sep.base east) {example.com};
    \node (pto) at (host.base east) {:443};
    \node (ruta) at (pto.base east) {/docs/?q=tls};
  \end{scope}
  \foreach \p in {esq,host,pto,ruta}
    \draw[linea,line width=1.2pt] ([xshift=1pt]\p.south west) -- ([xshift=-1pt]\p.south east);
  \node[cajag,anchor=north] (b1) at (2.0,-1.05)
    {\arriba{1.05cm}{3.5cm}{\textbf{Esquema}\\exige una conexión asegurada con TLS}};
  \node[cajag,anchor=north] (b2) at (6.47,-1.05)
    {\arriba{1.05cm}{3.5cm}{\textbf{Host}\\DNS lo traduce a una dirección IP}};
  \node[cajag,anchor=north] (b3) at (10.94,-1.05)
    {\arriba{1.05cm}{3.5cm}{\textbf{Puerto}\\443 si la URL no indica otro}};
  \node[cajaa,inner sep=6pt,font=\footnotesize,anchor=north] (b4) at (15.4,-1.05)
    {\arriba{1.05cm}{3.5cm}{\textbf{Ruta y consulta}\\van dentro de la petición HTTP, cifradas}};
  \draw[rel] (esq.south) -- ++(0,-0.3) -| (b1.north);
  \draw[rel] (host.south) -- ++(0,-0.3) -| (b2.north);
  \draw[rel] (pto.south) -- ++(0,-0.6) -| (b3.north);
  \draw[rel] (ruta.south) -- ++(0,-0.3) -| (b4.north);
\end{tikzpicture}
\end{bloque}

\textbf{Paso 1: DNS.} El resolvedor traduce el nombre a una IP consultando servidores de
nombres en el puerto 53 \fuente{RFC 1035, §4.2}. Puede necesitar varias consultas y un tiempo
arbitrario, y guarda las respuestas en una caché \fuente{RFC 1035, §2.2}: por eso este resumen
no lo cuenta entre los RTT. Una respuesta falsa puede llevar a otra IP; verificar el
certificado detecta a quien controla la resolución \fuente{RFC 9110, §4.3.4}.

\textbf{Paso 2: TCP.} El cliente envía un SYN, el servidor responde con SYN,ACK y, al
recibirlo, el cliente ya puede mandar datos junto con su ACK \fuente{RFC 9293, §3.5, fig. 6}.
La conexión cuesta un RTT y prepara el camino para TLS.

\ojo[Confusión típica]{Creer que llegar a la IP y al puerto 443 prueba que se habla con el
sitio. HTTPS no usa TCP ni el puerto para decidir la autoridad, porque quedan fuera de la
comunicación asegurada; la prueba es el certificado \fuente{RFC 9110, §4.3.3}.}

\begin{repaso}
\item ¿Qué parte de la URL viaja cifrada, y dónde viaja el host en claro?
\item ¿Por qué el DNS no entra en la cuenta de RTT?
\item ¿Por qué llegar a la IP y al puerto correctos no prueba la identidad del servidor?
\end{repaso}

% ============================ 02 ============================
\section{TLS 1.3 acuerda las claves en un solo viaje de ida y vuelta}\label{sec:handshake}

\textbf{Definición.} El handshake es «una negociación inicial entre cliente y servidor que
establece los parámetros de sus interacciones posteriores» \fuente{RFC 8446, §1.1}. En
TLS 1.3 tiene tres fases \fuente{RFC 8446, §2}:

\begin{itemize}
\item \textbf{Intercambio de claves.} El ClientHello ofrece versiones y cifrados y lleva su
  parte del acuerdo de claves (\cod{key_share}); el ServerHello responde con la suya. De ahí
  salen las claves, y lo que sigue va cifrado.
\item \textbf{Parámetros del servidor.} EncryptedExtensions, y CertificateRequest si pide
  certificado al cliente, algo raro en la web \fuente{MDN-TLS}.
\item \textbf{Autenticación.} Certificate lleva el certificado; CertificateVerify, una firma
  del handshake con la clave privada; Finished, un código de autenticación (MAC) del
  handshake que confirma las claves.
\end{itemize}

\begin{bloque}[\textbf{Ejemplo.} Una visita nueva: tres RTT hasta la primera respuesta \fuente{RFC 9293, fig. 6 · RFC 8446, fig. 1}. Sobre celeste, lo que viaja cifrado: todo lo que sigue al ServerHello, con claves del handshake \{…\} o de la aplicación {[\,…\,]}.]
\begin{tikzpicture}[msg/.style={-{Stealth[length=5pt]},line width=0.8pt,color=rotulo},
    cif/.style={-{Stealth[length=5pt]},line width=0.9pt,color=acento},
    txt/.style={font=\scriptsize,inner sep=2pt,sloped,above,text=rotulo},
    txtc/.style={txt,text=acento}]
  \path (0,0.6) (17.4,0);
  \def\xc{4.4} \def\xs{16.2} \def\d{1.2}
  \node[font=\footnotesize\bfseries,text=tinta] at (\xc,0.45) {Cliente};
  \node[font=\footnotesize\bfseries,text=tinta] at (\xs,0.45) {Servidor};
  \fill[suave] (\xs,-3*\d) -- (\xc,-4*\d) -- (\xc,-6*\d-0.15) -- (\xs,-6*\d-0.15) -- cycle;
  \draw[linea,line width=1.2pt] (\xc,0.2) -- (\xc,-6*\d-0.15);
  \draw[linea,line width=1.2pt] (\xs,0.2) -- (\xs,-6*\d-0.15);
  \draw[msg] (\xc,0) -- node[txt] {SYN} (\xs,-\d);
  \draw[msg] (\xs,-\d) -- node[txt] {SYN,ACK} (\xc,-2*\d);
  \draw[msg] (\xc,-2*\d) -- node[txt] {ACK · ClientHello + key\_share + server\_name} (\xs,-3*\d);
  \draw[msg] (\xs,-3*\d) -- node[txt] {ServerHello + key\_share}
    node[txtc,below,align=center] {\{EncryptedExtensions, Certificate,\\CertificateVerify, Finished\}} (\xc,-4*\d);
  \draw[cif] (\xc,-4*\d) -- node[txtc,below] {\{Finished\} · [GET /docs/?q=tls]} (\xs,-5*\d);
  \draw[cif] (\xs,-5*\d) -- node[txtc,below] {[200 OK y la página]} (\xc,-6*\d);
  \foreach \ya/\yb/\r in {0/2/TCP · 1 RTT, 2/4/TLS 1.3 · 1 RTT, 4/6/HTTP · 1 RTT} {
    \draw[rotulo,line width=0.6pt] (\xc-0.25,-\ya*\d-0.05) -- (\xc-0.45,-\ya*\d-0.05)
      -- (\xc-0.45,-\yb*\d+0.05) -- (\xc-0.25,-\yb*\d+0.05);
    \node[etiqueta,anchor=east,font=\scriptsize\bfseries] at (\xc-0.55,{-(\ya+\yb)/2*\d}) {\r};
  }
\end{tikzpicture}
\end{bloque}

\noindent\begin{columna}{0.55\linewidth}
\figura{fig/f1_viajes.pdf}{RTT hasta la primera respuesta, sin el DNS: cuenta propia sobre los diagramas de las RFC. En negrita, TLS 1.3 completo.}
\end{columna}\hfill
\begin{columna}{0.41\linewidth}
TLS 1.3 ahorra un RTT frente a TLS 1.2, que esperaba dos antes de los datos
\fuente{RFC 5246, fig. 1}. Si el servidor no acepta el \cod{key_share} del cliente, responde
con HelloRetryRequest, el ClientHello se repite y el ahorro se pierde
\fuente{RFC 8446, §2.1}.

Con una clave precompartida (PSK) de una conexión anterior, el cliente puede mandar la
petición junto al ClientHello: es el modo 0-RTT \fuente{RFC 8446, §2.2 y §2.3}.
\end{columna}

\ojo[Confusión típica]{Suponer que los datos 0-RTT tienen las garantías del resto. No tienen
secreto hacia adelante, porque se cifran solo con claves de la PSK, y nada impide repetirlos
entre conexiones \fuente{RFC 8446, §2.3}. Secreto hacia adelante: una clave de largo plazo
filtrada después no compromete la sesión \fuente{RFC 8446, §E.1}.}

\begin{repaso}
\item ¿Qué lleva el ClientHello para que las claves salgan en un solo RTT?
\item ¿Qué prueba CertificateVerify y qué confirma Finished?
\item ¿Cuándo se pierde el RTT que ahorra TLS 1.3?
\item ¿Qué dos garantías les faltan a los datos 0-RTT?
\end{repaso}

% ============================ 03 ============================
\section{El certificado ata un nombre a una clave, y la cadena lo avala}\label{sec:certificados}

\textbf{Definición.} Validar una ruta de certificación es verificar el vínculo entre el
nombre del sujeto y la clave pública del certificado final a partir de la clave de un ancla
de confianza \fuente{RFC 5280, §6.1}. El certificado lleva además el emisor, la vigencia y los
nombres DNS en \cod{subjectAltName} \fuente{RFC 5280, §4.1.2.5 y §4.2.1.6}; solo uno con
\cod{cA} verdadero en \cod{basicConstraints} firma otros \fuente{RFC 5280, §4.2.1.9}. El
servidor manda su certificado y después los que lo certifican; el ancla puede faltar, porque
el cliente ya la tiene \fuente{RFC 8446, §4.4.2}.

\begin{bloque}[\textbf{Ejemplo.} La cadena por defecto de un certificado RSA que emite Let's Encrypt YR1. Root YR aún no está en los almacenes de los navegadores: por eso la firma también ISRG Root X1, el ancla, en azul \fuente{LE-cadenas}.]
\begin{tikzpicture}
  \path (0,0.5) (17.4,0);
  \node[etiqueta,font=\scriptsize\ttfamily,text=rotulo,anchor=south west] at (0,0.62) {LO ENVÍA EL SERVIDOR};
  \node[etiqueta,font=\scriptsize\ttfamily,text=rotulo,anchor=south west] at (13.78,0.62) {LO TIENE EL NAVEGADOR};
  \node[cajag,anchor=north west] (ee) at (0,0.5)
    {\arriba{1.05cm}{3.2cm}{\textbf{example.com}\\certificado final: nombre y clave pública}};
  \node[cajag,anchor=north west] (yr1) at (4.58,0.5)
    {\arriba{1.05cm}{3.2cm}{\textbf{Let's Encrypt YR1}\\intermedia RSA 2048, vigente hasta 2028-09-02}};
  \node[cajag,anchor=north west] (ryr) at (9.16,0.5)
    {\arriba{1.05cm}{3.2cm}{\textbf{Root YR}\\raíz nueva, fuera de los almacenes}};
  \node[cajaa,inner sep=6pt,font=\footnotesize,anchor=north west] (x1) at (13.78,0.5)
    {\arriba{1.05cm}{3.2cm}{\textbf{ISRG Root X1}\\confiable hasta 2030-06-04}};
  \draw[linea,dashed,line width=0.9pt] (13.47,1.0) -- (13.47,-1.35);
  \draw[rel] (yr1.west) -- node[relacion,above] {firma} (ee.east);
  \draw[rel] (ryr.west) -- node[relacion,above] {firma} (yr1.east);
  \draw[rel] (x1.west) -- node[relacion,above,pos=0.3] {firma} (ryr.east);
\end{tikzpicture}
\end{bloque}

De cada certificado, el cliente verifica la firma, la vigencia, que no esté revocado y que
su emisor sea el sujeto del anterior \fuente{RFC 5280, §6.1.3}. Después compara el nombre de la
URL con el certificado final.

\begin{bloque}[Cualquier «no» lleva al error, con borde terracota: el navegador tiene que pedir confirmación al usuario o cortar la conexión \fuente{RFC 9110, §4.3.4}.]
\begin{tikzpicture}
  \path (0,0) (17.4,0);
  \node[cajag,anchor=north west] (q1) at (0,0)
    {\arriba{0.65cm}{3.4cm}{\textbf{¿Cada certificado pasa los cuatro controles?}}};
  \node[cajag,anchor=north west] (q2) at (4.52,0)
    {\arriba{0.65cm}{3.4cm}{\textbf{¿La cadena llega a una raíz del almacén?}}};
  \node[cajag,anchor=north west] (q3) at (9.04,0)
    {\arriba{0.65cm}{3.4cm}{\textbf{¿El nombre de la URL está en el certificado?}}};
  \node[cajag,anchor=north west] (ok) at (13.56,0)
    {\arriba{0.65cm}{3.4cm}{\textbf{Identidad verificada}\\sigue el handshake}};
  \draw[flecha] (q1) -- node[relacion,above] {sí} (q2);
  \draw[flecha] (q2) -- node[relacion,above] {sí} (q3);
  \draw[flecha] (q3) -- node[relacion,above] {sí} (ok);
  \coordinate (bus) at (0,-1.4);
  \foreach \q in {q1,q2,q3} {
    \draw[bus] (\q.south) -- (\q.south |- bus);
    \node[relacion,anchor=west] at ($(\q.south)+(0.05,-0.16)$) {no};
  }
  \draw[bus] (q1.south |- bus) -- (q3.south |- bus);
  \node[cajag,fill=suave2,draw=acento2,line width=0.8pt,anchor=north] (err) at ($(q2.south |- bus)+(0,-0.35)$)
    {\textbf{Error de certificado:} confirmar con el usuario o cortar};
  \draw[rel] (q2.south |- bus) -- (err.north);
\end{tikzpicture}
\end{bloque}

\ojo[Confusión típica]{Buscar el nombre del sitio en el Common Name (CN). El cliente arma
la identidad de referencia como DNS-ID con el host de la URL; una CN-ID no debe usarse
\fuente{RFC 9110, §4.3.4}.}

\begin{repaso}
\item ¿Qué une un certificado X.509, y qué campo dice si puede firmar otros?
\item ¿Qué cuatro controles pasa cada certificado de la ruta?
\item ¿Por qué Root YR necesita la firma de ISRG Root X1?
\item ¿Qué tiene que hacer el navegador si el nombre no coincide?
\end{repaso}

% ============================ 04 ============================
\section{HTTPS protege el contenido y la identidad, no los metadatos}\label{sec:alcance}

Handshake y certificado determinan juntos lo que HTTPS protege: el servidor queda
autenticado y los datos, cifrados e íntegros \fuente{RFC 8446, §1}. Lo que viaja fuera de TLS
lo ve cualquiera en la red.

\begin{bloque}[A la izquierda, lo que ve la red aunque se use HTTPS \fuente{RFC 9110, §4.3.3 · RFC 8446, §E.3}; a la derecha, lo que queda dentro de TLS \fuente{MDN-TLS}. En terracota, el nombre en el ClientHello: va en claro porque solo se cifra lo que sigue al ServerHello \fuente{RFC 8446, §1.2}.]
\begin{tikzpicture}
  \node[cajag,anchor=north west] (red) at (0,0)
    {\arriba{1.1cm}{8.1cm}{\familia{Lo ve la red}{IP y puerto, nombre en la consulta DNS, {\color{acento2}\bfseries nombre en el ClientHello}, tamaño y ritmo del tráfico}}};
  \node[cajag,anchor=north west] (tls) at (8.95,0)
    {\arriba{1.1cm}{8.03cm}{\familia{Solo lo ven cliente y servidor}{ruta y consulta, encabezados, cookies, cuerpo de la petición, cuerpo de la respuesta}}};
  \draw[acento2,dashed,line width=0.9pt] (8.72,0.1) -- (8.72,-1.65);
  \node[fill=white,inner sep=1.5pt,rotate=90,font=\scriptsize\bfseries,text=acento2] at (8.72,-0.77) {TLS};
\end{tikzpicture}
\end{bloque}

\begin{bloque}[HTTPS impide o detecta los ataques al contenido y a la identidad; la primera visita, el 0-RTT y los metadatos necesitan otras defensas.]
\small
\begin{tabularx}{\linewidth}{@{}P{4.6cm}LP{2.6cm}@{}}
\toprule
\enc{Amenaza} & \enc{Qué hace HTTPS} & \enc{Fuente} \\
\midrule
Alguien en el medio lee o cambia el tráfico & Lo impide: cifrado, integridad y autenticación del servidor & \fuente{MDN-TLS} \\
Una respuesta DNS falsa lleva a otra IP & Lo detecta: el certificado no coincide con el nombre & \fuente{RFC 9110, §4.3.4} \\
Degradar a http la primera visita (SSL stripping) & HSTS lo evita en las siguientes; la lista de precarga, también en la primera & \fuente{MDN-TLS} \\
Un script pedido por http en una página https (contenido mixto) & El navegador pasa la petición a https o la bloquea & \fuente{MDN-TLS} \\
Repetir datos 0-RTT & TLS no lo impide; el servidor puede limitarlo & \fuente{RFC 8446, §8} \\
Deducir el contenido por el tamaño y el ritmo & No lo impide; el relleno de los registros lo dificulta & \fuente{RFC 8446, §E.3} \\
\bottomrule
\end{tabularx}
\end{bloque}

HSTS es el encabezado \cod{Strict-Transport-Security}: pide al navegador que las próximas
visitas vayan directo por https \fuente{MDN-TLS}. Con DoH, la consulta DNS viaja dentro de
HTTPS y los equipos del camino no pueden interferir \fuente{RFC 8484, §1}.

\begin{repaso}
\item ¿Qué datos de una visita ve un observador de la red, y por qué?
\item ¿Por qué una respuesta DNS falsa no alcanza para suplantar un sitio https?
\item ¿Qué evita HSTS, y qué agrega la lista de precarga?
\end{repaso}

\anexo{Referencias}
\begin{referencias}
\referencia{RFC 8446}{E. Rescorla, \emph{The Transport Layer Security (TLS) Protocol Version 1.3}, IETF, 2018: \url{https://www.rfc-editor.org/rfc/rfc8446}.}
\referencia{RFC 5246}{T. Dierks y E. Rescorla, \emph{The TLS Protocol Version 1.2}, IETF, 2008: \url{https://www.rfc-editor.org/rfc/rfc5246}.}
\referencia{RFC 9110}{R. Fielding, M. Nottingham y J. Reschke (eds.), \emph{HTTP Semantics}, IETF, 2022: \url{https://www.rfc-editor.org/rfc/rfc9110}.}
\referencia{RFC 9293}{W. Eddy (ed.), \emph{Transmission Control Protocol (TCP)}, IETF, 2022: \url{https://www.rfc-editor.org/rfc/rfc9293}.}
\referencia{RFC 1035}{P. Mockapetris, \emph{Domain Names — Implementation and Specification}, IETF, 1987: \url{https://www.rfc-editor.org/rfc/rfc1035}.}
\referencia{RFC 8484}{P. Hoffman y P. McManus, \emph{DNS Queries over HTTPS (DoH)}, IETF, 2018: \url{https://www.rfc-editor.org/rfc/rfc8484}.}
\referencia{RFC 5280}{D. Cooper y otros, \emph{Internet X.509 Public Key Infrastructure Certificate and CRL Profile}, IETF, 2008: \url{https://www.rfc-editor.org/rfc/rfc5280}.}
\referencia{MDN-TLS}{MDN Web Docs, \emph{Transport Layer Security (TLS)}: \url{https://developer.mozilla.org/en-US/docs/Web/Security/Defenses/Transport_Layer_Security}.}
\referencia{LE-cadenas}{Let's Encrypt, \emph{Chains of Trust}, actualizada el 8 de julio de 2026: \url{https://letsencrypt.org/certificates/}.}
\referencia{LE-validación}{Let's Encrypt, \emph{How It Works}: \url{https://letsencrypt.org/how-it-works/}.}
\end{referencias}

\end{document}
